\documentclass[../../../include/open-logic-section]{subfiles}

\begin{document}

\olfileid{sth}{replacement}{strength}
\olsection{Kakuwataning Panggantian}

Kita miwiti kanthi pangamatan prasaja babagan
kakuwataning Panggantian: yen ora ngluwihi $\Z$,
kita ora bisa mbuktekake anane ordinal von Neumann
sebarang kang luwih gedhe utawa padha karo
$\omega + \omega$.

Mangkene garis gedhe alesane. Makarya ing~$\ZF$,
gatekna himpunan $V_{\omega+\omega}$. Himpunan
iki dadi domain kanggo sawijining \emph{model}
saka~$\Z$. Kanggo ndeleng iki, kita ngenalake
notasi kanggo \emph{relativisasi} formula:
\begin{defn}\ollabel{formularelativization}
	Kanggo saben himpunan $M$ lan formula $\phi$,
	$\phi^M$ yaiku formula kang dipikolehi kanthi
	matesi kabeh kuantor ing $\phi$ marang $M$.
	Tegese, ganti ``$\lexists[x]$'' dadi
	``$(\lexists[x \in M])$'', lan ganti
	``$\lforall[x]$'' dadi
	``$(\lforall[x \in M])$''.
\end{defn}\noindent
Bisa dibuktekake yen kanggo saben aksioma
$\phi$ saka~$\Z$, kita duwe
$\ZF \vdash \phi^{V_{\omega+\omega}}$.
Nanging $\omega+\omega$ ora \emph{kalebu ing}
$V_{\omega+\omega}$ miturut
\olref[spine][rank]{ordsetrankalpha}. Dadi $\Z$
konsisten karo ora anane $\omega+\omega$.

Iki sebabe kita kandha ing
\olref[ordinals][replacement]{sec} yen
\olref[ordinals][ordtype]{thmOrdinalRepresentation}
ora bisa dibuktekake tanpa Panggantian. Awit
ing~$\Z$ gampang netepake urutan becik kang
eksplisit, kang miturut intuisi \emph{mesthine}
nduweni tipe urutan $\omega+\omega$. Conto
ora formale wis diwenehake ing
\olref[ordinals][idea]{sec}, nalika kita
nerangake urutan wilangan asli iki:
\begin{align*}
	n \lessdot m \text{ yen lan mung yen }&\text{salah siji }n < m\text{ lan }m-n\text{ genep,}\\
	& \text{utawa $n$ genep lan $m$ ganjil.}
\end{align*}
Nanging yen $\omega+\omega$ ora ana, urutan
becik iki ora isomorfis karo ordinal apa wae.
Mula $\Z$ \emph{ora} mbuktekake
\olref[ordinals][ordtype]{thmOrdinalRepresentation}.

Kosokbaline, Panggantian ngidini kita
mbuktekake anane $\omega+\omega$, lan mula
uga anane $V_{\omega+\omega}$. Ora mung kuwi.
Kanggo \emph{saben} urutan becik kang bisa
kita tetepake,
\olref[ordinals][ordtype]{thmOrdinalRepresentation}
ngandhakake yen ana sawijining $\alpha$ kang
isomorfis karo urutan becik iku, lan mula
$V_\alpha$ ana. Dadi Panggantian njamin
yen hierarki himpunan iku \emph{dhuwur banget}.

Ing sawetara bagean sabanjure, lan maneh ing
\olref[card-arithmetic][fix]{sec}, kita bakal
luwih ngerti \emph{sepira} dhuwure hierarki
kang diwajibake dening Panggantian. Saiki,
pokok prasajane yaiku Panggantian pancen
\emph{mbutuhake} pambeneran!

\end{document}
